\section*{Chapitre \ref{ch:hf:queue-reactive-models-and-large-tick-assets} --- Modèles réactifs à la file et actifs à gros tick}

\begin{solution}{exo:hf:queue-reactive-models-and-large-tick-assets:1}
$290/200=1.45$ lots par seconde.
\end{solution}

\begin{solution}{exo:hf:queue-reactive-models-and-large-tick-assets:2}
$s^b=3/15=0.2$~: inférieur à 0.3, donc il ne rejoint pas le cours acheteur~; la part du cours vendeur vaut 0.8, donc il rejoint le cours vendeur.
\end{solution}

\begin{solution}{exo:hf:queue-reactive-models-and-large-tick-assets:3}
$0.392\times3\,400\times\$0.01=\$13.33$.
\end{solution}

\begin{solution}{exo:hf:queue-reactive-models-and-large-tick-assets:4}
Une file opposée courte se vide vite, et le prix s'éloigne alors de l'ordre au repos avant que celui-ci ne soit atteint~: l'ordre reste non exécuté.
Avec une file opposée longue, le prix risque peu de s'éloigner, si bien que même le fond de la file finit par être exécuté.
\end{solution}

\begin{solution}{exo:hf:queue-reactive-models-and-large-tick-assets:5}
Rejoindre sur une part favorable ne sélectionne les bons états qu'au moment de l'entrée. Les ordres qui restent sont alors exécutés surtout après que la part a
tourné contre eux, alors que leur file se consomme devant un mouvement de prix~: la règle ne garde précisément que les ordres dont les exécutions sont défavorables
($-0.117$ ticks par action, contre 0.026 en rejoignant systématiquement).
\end{solution}

\begin{solution}{exo:hf:queue-reactive-models-and-large-tick-assets:6}
Un seuil plus élevé ne retient que les meilleurs états, si bien que chaque exécution est meilleure, mais il cote moins souvent et exécute moins d'actions~; l'edge par session est
le produit, maximal en 0.3 (13.33~\$), tandis que le mark-out par action est maximal en 0.4 (0.484).
\end{solution}

\begin{solution}{exo:hf:queue-reactive-models-and-large-tick-assets:7}
Probabilité d'exécution de 0.825, valeur de 0.594 tick par action, mouvement adverse $-0.22$~: après ses exécutions, le mid se déplace en sa faveur en moyenne, parce que la
file vendeuse courte tend à se vider et le prix à monter.
\end{solution}

\begin{solution}{exo:hf:queue-reactive-models-and-large-tick-assets:8}
Le modèle ne comporte pas de flux informé~: sa seule sélection adverse est mécanique. Les traders informés du marché simulé dégradent les exécutions précisément lorsque
les files se consomment, ce que les valeurs du modèle n'intègrent pas~; le teneur de marché du chapitre~1, toujours présent dans les files, a perdu sa fourchette. Mesurer les mark-out de
ses propres exécutions et ne coter que là où ils sont positifs.
\end{solution}

\begin{solution}{pb:hf:queue-reactive-models-and-large-tick-assets:1}
\begin{enumerate}
\item La fourchette vaut presque toujours un tick lorsque les meilleures files sont longues~; le teneur de marché contrôle les files où il se place et le moment où il les quitte, non le prix.
\item Un tick dans 96.3~\% du temps~; médiane de la meilleure file~: 17 lots.
\item Les ordres au marché sont ici trop petits pour épuiser une meilleure file, si bien que le prix suivant n'est presque jamais atteint.
\item Voir la \cref{def:hf:queue-reactive-models-and-large-tick-assets:qr}.
\item La \Cref{met:hf:queue-reactive-models-and-large-tick-assets:estimate}~; les ordres à cours limité entre 1.3 et 2.3 lots par seconde, sans tendance, les annulations passant
de 0.34 à environ 1.5, les ordres au marché tombant de 0.55 à environ 0.3.
\item Chaque ordre au repos s'annule à son propre rythme, si bien qu'une file plus longue perd ses lots plus vite.
\item Un lot dans la file acheteuse, avec $n$ devant lui~; entrées, annulations et ordres au marché aux taux ajustés~; exécuté lorsqu'il est atteint, abandonné si la file vendeuse
se vide ou après 60 secondes~; après une exécution, évalué au mid dix secondes plus tard~: $0.5$ tick plus le mouvement du mid.
\item 0.60 (tête de file) et 0.32 (fond de file) à 10 contre 2~; 0.12 pour les deux à 2 contre 20.
\item La file est sur le point d'être consommée~: chaque place est exécutée, suivie du mouvement de prix.
\item Le modèle ne comporte pas de flux informé~; l'information contenue dans le flux est omise.
\item Voir la \cref{def:hf:queue-reactive-models-and-large-tick-assets:rules}.
\item Rejoindre toujours 0.026, rejoindre et rester $-0.117$, rejoindre et quitter 0.392 ticks par action~; dans le modèle, la tête d'une file de dix lots face
à un cours vendeur de deux lots, elle vaut 0.60 ticks et le fond 0.32.
\item Il évalue l'inventaire de fin de chaque session, dont les variations sont grandes face aux écarts des règles sur six sessions.
\item Lorsque l'autre côté est sur le point de consommer la file, la tête est exécutée en premier, et le mouvement suit.
\item 0.3 pour l'edge par session~; le mark-out continue de monter jusqu'à 0.4, car des règles plus strictes retiennent de meilleures exécutions, mais en moins grand nombre.
\item Toujours 365, rester 245, quitter 468, quitter sauf en tête 448~: quitter signifie annuler puis rejoindre la file au fil des mouvements du carnet d'ordres.
\item Ajuster le mark-out des exécutions, par état, aux variables du carnet d'ordres au moment de l'exécution et l'ajouter au gain, ou ajouter des ordres au marché informés dont le taux
dépend d'un prix efficient caché.
\item Avec une allocation au prorata, la position dans la file ne compte pas, c'est la taille qui compte (chapitre~15)~; quitter ne coûte aucune place.
\item La fourchette vaut plusieurs ticks, si bien que le choix du prix réapparaît (chapitres~3--4) et que les files sont courtes~; les règles comptent moins.
\item Bien quitter.
\end{enumerate}
\end{solution}

\begin{solution}{iq:hf:queue-reactive-models-and-large-tick-assets:1}
La file vendeuse est sur le point de se vider et le prix de monter~: un ordre d'achat en tête de file sera probablement bien exécuté, et le mouvement suit en sa faveur. Il faut rester~; le
danger est l'inverse~: une file acheteuse courte.
\iqlookfor{lecture de la file du côté opposé, et le fait que la tête de file n'est pas toujours sûre.}
\end{solution}

\begin{solution}{iq:hf:queue-reactive-models-and-large-tick-assets:2}
Enregistrer la position de chaque ordre à chaque événement et son issue (exécution, mark-out à un horizon donné ou annulation)~; estimer le mark-out espéré et la probabilité
d'exécution par place et par taille des files, hors échantillon~; l'écart entre places est la valeur de la position dans la file.
\iqlookfor{issues conditionnelles à la place, et biais de sélection des ordres qui restent.}
\end{solution}

\begin{solution}{iq:hf:queue-reactive-models-and-large-tick-assets:3}
Un ordre situé au fond de la file n'est atteint que lorsque l'essentiel de la file a été consommé, c'est-à-dire quand le prix est sur le point de le traverser~: une probabilité d'exécution plus élevée
tardive dans la vie d'une file s'accompagne de moins bonnes exécutions.
\iqlookfor{le caractère conditionnel des exécutions tardives.}
\end{solution}

\begin{solution}{iq:hf:queue-reactive-models-and-large-tick-assets:4}
À partir d'un flux ordre par ordre~: à son accusé de réception, l'ordre se place derrière chaque ordre au repos à son prix~; chaque exécution ou annulation
d'un ordre devant lui le fait avancer. Les ordres cachés, les paquets perdus ou réordonnés et les règles de priorité propres aux plateformes (changements de taille, remplacements) faussent le décompte.
\iqlookfor{suivi ordre par ordre et ses modes de défaillance.}
\end{solution}

\begin{solution}{iq:hf:queue-reactive-models-and-large-tick-assets:5}
Chaque annulation suivie d'une nouvelle entrée renvoie l'ordre au fond de la file~: sur une action à un tick, la place fait l'essentiel de la valeur de l'ordre, si bien que ce va-et-vient la
lui fait perdre, en plus d'ajouter des messages.
\iqlookfor{priorité perdue lors d'un remplacement.}
\end{solution}

\begin{solution}{iq:hf:queue-reactive-models-and-large-tick-assets:6}
Avec $k$ ordres devant lui, les événements parmi eux surviennent au taux $\mu+k\nu$~; le prochain est un ordre au marché avec la probabilité $\mu/(\mu+n\nu)$. Chaque événement retire un
ordre devant lui, si bien que le temps espéré jusqu'à ce qu'il n'en reste aucun est $\sum_{k=1}^n1/(\mu+k\nu)$.
\iqlookfor{horloges exponentielles concurrentes et somme des étapes.}
\end{solution}
